\PassOptionsToPackage{table}{xcolor}
\documentclass[10pt, aspectratio=169, squeeze]{beamer}

% --- Tema e Cores ---
\usetheme{metropolis} % Tema moderno e limpo
\metroset{block=fill} % Blocos com preenchimento de cor para destaque

% --- Pacotes ---
\usepackage[utf8]{inputenc}
\usepackage[T1]{fontenc}
\usepackage[brazil]{babel}
\usepackage{booktabs}
\usepackage{tabularx}
\usepackage{colortbl}
\usepackage{tikz}

% --- Definição de Cores Personalizadas ---
\definecolor{primaryColor}{RGB}{0, 51, 102}    % Azul Marinho (Base)
\definecolor{accentColor}{RGB}{178, 34, 34}    % Vermelho (Dívidas)
\definecolor{successColor}{RGB}{34, 139, 34}   % Verde (Acordo)
\definecolor{lightGray}{RGB}{245, 245, 245}
\definecolor{secondaryText}{RGB}{100, 100, 100} % Cinza para textos secundários

% --- Configuração do Tema ---
\setbeamercolor{palette primary}{bg=primaryColor, fg=white}
\setbeamercolor{alerted text}{fg=accentColor}
\setbeamercolor{example text}{fg=successColor}
\setbeamercolor{block title}{bg=primaryColor, fg=white}
\setbeamercolor{block title alerted}{bg=accentColor, fg=white}

% --- Metadados ---
\title{Plano de Dissolução e Partilha}
\subtitle{Relatório Técnico e Proposta Financeira}
\author{\textbf{Studio Dental Comércio de Higiene Oral Ltda}}
\date{Data Base: 25 de Novembro de 2025}
\titlegraphic{\hfill\scriptsize CNPJ: 46.125.234/0001-66}

% --- Macro para Moeda ---
\newcommand{\reais}[1]{R\$ #1}

\begin{document}

% --- Slide de Capa ---
\maketitle

% --- Slide 1: Visão Geral ---
\begin{frame}{Quadro Societário Atual}
    \begin{block}{Objetivo}
        Dissolução societária com partilha de bens e reconhecimento de passivos.
    \end{block}
    
    \vspace{0.5cm}
    
    \begin{alertblock}{Cenário Base}
        Reconhecimento de passivo de locação retroativo e inviabilidade de venda externa dos ativos.
    \end{alertblock}
\end{frame}

% --- Slide 1b: Distribuição de Cotas ---
\begin{frame}{Distribuição de Cotas}
    \centering
    \begin{table}
        \rowcolors{1}{lightGray}{white}
        \begin{tabular}{l c}
            \toprule
            \textbf{Sócio} & \textbf{\%} \\
            \midrule
            Frederique A. Lisboa & 40\% \\
            Kroner M. Costa & 40\% \\
            Eduardo B. Nunes & 10\% \\
            Alessandra S. Alves & 5\% \\
            Simone H. dos Santos & 5\% \\
            \bottomrule
        \end{tabular}
    \end{table}
\end{frame}

% --- Slide 2: Premissas ---
\begin{frame}{Premissas Fundamentais do Acordo}
    Para viabilizar uma saída justa, estabelecem-se os seguintes pilares:
    \vspace{0.5cm}
    \begin{enumerate}
        \item \textbf{Passivo de Locação:} Reconhecimento da dívida pelo uso do imóvel (propriedade do sócio Kroner) desde Jan/2014.
        \item \textbf{Dação em Pagamento:} Kroner aceita ativos fixos e infraestrutura como abatimento parcial da dívida.
        \item \textbf{Liberação de Ativos Móveis:} Os demais sócios retiram seus equipamentos de trabalho sem ônus.
        \item \textbf{Passivo Bancário:} Continua sendo rateado proporcionalmente até a liquidação (2031).
    \end{enumerate}
\end{frame}

% --- Slide 3: As Dívidas (O Problema) ---
\begin{frame}{Análise dos Passivos (Dívidas)}
    \begin{columns}[T]
        \column{0.48\textwidth}
            \begin{alertblock}{1. Passivo de Locação (Retroativo)}
                Dívida com o sócio \textbf{Kroner Machado}.
                \begin{itemize}
                    \item Período: 142 meses (11 anos).
                    \item Valor Mensal Base: \reais{10.000,00}.
                \end{itemize}
                \vspace{0.2cm}
                \centering \textbf{Total: \reais{1.420.000,00}}
            \end{alertblock}

        \column{0.48\textwidth}
            \begin{block}{2. Passivo Bancário (Mensal)}
                Dívida Externa (Banco).
                \begin{itemize}
                    \item Parcela Mensal: \reais{7.300,00}.
                    \item Rateio: Proporcional às cotas.
                \end{itemize}
                \vspace{0.2cm}
                \centering \small (Frederique/Kroner: \reais{2.920} cada)
            \end{block}
    \end{columns}
\end{frame}

% --- Slide 4: Partilha - Grupo A ---
\begin{frame}{Partilha: Grupo A (Frederique e Eduardo)}
    \textit{Laboratório de Prótese Completo + Tecnologia CAD/CAM.}
    
    \begin{table}
        \centering
        \small
        \begin{tabularx}{\textwidth}{X r}
            \toprule
            \textbf{Principais Itens} & \textbf{Valor Real (Depreciado)} \\
            \midrule
            Ceramill Motion 2 (x2 Unidades) & 111.888,00 \\
            Scanner 3Shape (Unid. 1) & 50.000,00 \\
            Fornos, Impressoras 3D e Softwares & 47.500,00 \\
            Outros Equipamentos e Mobiliário & 71.862,00 \\
            \midrule
            \rowcolor{primaryColor!10} \textbf{TOTAL GRUPO A} & \textbf{\reais{281.250,00}} \\
            \bottomrule
        \end{tabularx}
    \end{table}
\end{frame}

% --- Slide 5: Partilha - Grupos B e C ---
\begin{frame}{Partilha: Clínico (Alessandra e Simone)}
    \begin{columns}
        \column{0.5\textwidth}
            \textbf{GRUPO B (Alessandra)} \\
            \textit{Equipamentos Portáteis e TV}
            \vspace{0.2cm}
            \begin{itemize}
                \item Kit Kavo, TV Samsung.
                \item Instrumental Diverso.
            \end{itemize}
            \vspace{0.2cm}
            \colorbox{lightGray}{\textbf{Total: \reais{28.010,00}}}

        \column{0.5\textwidth}
            \textbf{GRUPO C (Simone)} \\
            \textit{Kits Clínicos e TI}
            \vspace{0.2cm}
            \begin{itemize}
                \item iPad Pro, Câmera Nikon.
                \item Kit Kavo, PC Lenovo.
            \end{itemize}
            \vspace{0.2cm}
            \colorbox{lightGray}{\textbf{Total: \reais{27.410,00}}}
    \end{columns}
\end{frame}

% --- Slide 6: Partilha - Grupo D (Kroner) ---
\begin{frame}{Partilha: Grupo D (Kroner Machado)}
    \textit{Itens recebidos como \textbf{abatimento da dívida de aluguel}.}
    
    \begin{table}
        \centering
        \small
        \begin{tabularx}{\textwidth}{X r}
            \toprule
            \textbf{Categoria} & \textbf{Valor Real} \\
            \midrule
            \textbf{Infraestrutura:} Ar-condicionados, Móveis Planejados & 57.300,00 \\
            \textbf{Equipamentos:} 3 Cadeiras, Raio-X, Laser, Scanner & 219.250,00 \\
            \textbf{Intangível:} Goodwill e Ponto Comercial & 364.280,00 \\
            \midrule
            \rowcolor{primaryColor!10} \textbf{TOTAL GRUPO D} & \textbf{\reais{640.830,00}} \\
            \bottomrule
        \end{tabularx}
    \end{table}
\end{frame}

% --- Slide 7: Encontro de Contas (GRANDE FINAL) ---
\begin{frame}{Encontro de Contas e Conclusão}
    \small
    A matemática da dissolução demonstra que o acordo favorece a liberação dos sócios minoritários.

    \vspace{0.3cm}
    \begin{center}
    \begin{tikzpicture}[scale=0.85, transform shape]
        % Caixa Dívida
        \node[draw=accentColor, fill=accentColor!10, line width=1pt, rounded corners, align=center, minimum width=3.5cm, minimum height=1.2cm] (divida) at (0,0) {\textbf{Dívida Aluguel} \\ \reais{1.420.000}};
        
        % Sinal de Menos
        \node[font=\Large, text=secondaryText] at (2,0) {$-$};
        
        % Caixa Bens
        \node[draw=primaryColor, fill=primaryColor!10, line width=1pt, rounded corners, align=center, minimum width=3.5cm, minimum height=1.2cm] (bens) at (4.5,0) {\textbf{Bens Entregues} \\ \reais{640.830}};
        
        % Sinal de Igual
        \node[font=\Large, text=secondaryText] at (7,0) {$=$};
        
        % Saldo Restante
        \node[draw=black, fill=lightGray, dashed, align=center, minimum width=3.5cm, minimum height=1.2cm] (saldo) at (9.5,0) {\textbf{Saldo Restante} \\ \reais{779.170}};
        
        % Setas e Explicação
        \node[text=successColor, font=\small\bfseries, align=center] at (9.5, -1.2) {Valor Perdoado \\ por Kroner};
        \draw[->, thick, successColor] (9.5, -0.7) -- (9.5, -1.0);
    \end{tikzpicture}
    \end{center}

    \vspace{0.3cm}
    \textbf{Conclusão:} Kroner renuncia a \textbf{55\% da dívida} para permitir que os sócios levem seus laboratórios e kits de trabalho.
\end{frame}

% --- Slide 8: Assinaturas ---
\begin{frame}{Formalização}
    \centering
    \Large
    \textbf{Próximos Passos}
    \vspace{0.5cm}
    
    \normalsize
    Assinatura do distrato e efetivação da retirada dos bens móveis.
    
    \vspace{1cm}
    \textit{Goiânia, 25 de Novembro de 2025.}
\end{frame}

\end{document}
